% =============================================================
% CAPÍTULO 5 — Resultados de interpretabilidad mecánica
% Versión: matrícula de honor
% Alimentado por: NB04 (probing), NB05 (activation patching),
%                 NB06 (head ablation), NB07 (neuron analysis),
%                 NB08 (logit lens)
% =============================================================

\section{Resultados de interpretabilidad mecánica}
\label{sec:resultados-interpretabilidad}

El capítulo anterior documentó qué ocurre cuando se comprime BERT:
una transición de fase abrupta entre rango 384 y 256, una asimetría
radical entre componentes (factor 14$\times$ a rango~128 y hasta
72$\times$ a rango~64, normalizado por parámetros eliminados) y un
gradiente de fragilidad creciente con la profundidad. Este capítulo
aborda la pregunta complementaria: por qué ocurre. Para ello se
aplican cinco técnicas de interpretabilidad mecánica que examinan la
estructura interna del modelo desde ángulos diferentes pero convergentes.

La primera técnica, \emph{probing} (\S\ref{sec:probing}), entrena
clasificadores lineales sobre las representaciones intermedias para
determinar en qué capa se forma la señal emocional. La segunda,
\emph{logit lens} (\S\ref{sec:logit-lens}), aplica la cabeza de
clasificación entrenada del modelo a las representaciones de capas
intermedias para distinguir información disponible de información
alineada con la base del clasificador. La tercera, \emph{activation
patching} (\S\ref{sec:activation-patching}), restaura selectivamente
componentes del modelo original dentro de una versión comprimida para
localizar los loci funcionales críticos. La cuarta, \emph{ablación de
cabezas de atención} (\S\ref{sec:head-ablation}), anula individualmente
cada una de las 144 cabezas para construir una taxonomía funcional del
mecanismo atencional. La quinta, \emph{análisis neuronal de las FFN}
(\S\ref{sec:neuron-analysis}), examina la especialización de las 36\,864
neuronas de las redes feed-forward y su relación con la vulnerabilidad a
la compresión.

Cada sección presenta sus resultados de forma autónoma, pero el lector
observará un patrón transversal: las cinco técnicas convergen en una
arquitectura funcional de tres niveles, donde las capas tempranas extraen
señal léxica bruta, las capas intermedias realizan cómputo en una
geometría interna ilegible para el clasificador, y las capas profundas
(8--11) proyectan la representación sobre la base del clasificador. La
síntesis formal de esta convergencia se desarrolla en
\S\ref{sec:interp-sintesis}.


% =============================================================
% 5.1 — PROBING: CRISTALIZACIÓN EMOCIONAL POR CAPAS
% =============================================================

\subsection{Probing: cristalización emocional por capas}
\label{sec:probing}

\subsubsection{Diseño experimental}

Para determinar en qué punto de la jerarquía de capas de BERT emerge la
información emocional, se entrenan clasificadores \emph{probe} sobre las
representaciones intermedias del modelo. Se utiliza regresión logística
con regularización por defecto ($C = 1{,}0$) como clasificador, siguiendo
la metodología de Hewitt y Liang~\cite{hewitt2019designing}: el objetivo
del probe es medir la calidad de la representación, no la capacidad del
clasificador, por lo que se evita deliberadamente el tuning de
hiperparámetros. Un probe no-lineal complejo podría extraer información
de cualquier representación; un probe lineal solo detecta información ya
separada linealmente en el espacio de activaciones.

La representación de entrada es el token \texttt{[CLS]} extraído de cada
uno de los 13 \textit{hidden states} del modelo (la capa de embedding más
las 12 capas del encoder). Se entrena un probe binario independiente para
cada combinación de emoción y capa con una capa de información
significativa, resultando en 276 clasificadores activos (las 12 capas del
encoder por las 23 emociones; los probes sobre embedding producen
F1 = 0 uniformemente y se omiten del análisis efectivo).

Se define la \emph{capa de cristalización} de una emoción como la primera
capa en la que el F1 del probe alcanza el 80\% de su valor máximo a lo
largo de todas las capas. Este umbral es convencional; análisis de
sensibilidad con umbrales de 0{,}7 y 0{,}9 confirman que el orden
relativo de cristalización entre emociones permanece estable.

\subsubsection{Tres regímenes de cristalización}

Las 23 emociones se distribuyen en tres grupos claramente diferenciados
según su profundidad de cristalización
(Tabla~\ref{tab:cristalizacion-grupos},
Figura~\ref{fig:probing-heatmap}).

\begin{table}[ht]
\centering
\caption{Grupos de cristalización emocional. La capa indica la primera
capa donde el F1 del probe alcanza el 80\% de su máximo. La columna
``F1 en cap.'' muestra el F1 del probe en la capa de cristalización.}
\label{tab:cristalizacion-grupos}
\small
\begin{tabular}{llcl}
\toprule
\textbf{Grupo} & \textbf{Capas} & \textbf{N} & \textbf{Emociones representativas} \\
\midrule
Temprano & L0--L3 & 9 & gratitude (L0, 0{,}922), amusement (L0, 0{,}827), \\
         &        &   & love (L0, 0{,}818), curiosity (L0, 0{,}674), \\
         &        &   & optimism (L0, 0{,}569), admiration (L2, 0{,}735), \\
         &        &   & disgust (L2, 0{,}532), fear (L3, 0{,}708), \\
         &        &   & sadness (L3, 0{,}621) \\
\addlinespace
Medio    & L4--L7 & 8 & anger (L4, 0{,}543), remorse (L5, 0{,}655), \\
         &        &   & embarrassment (L5, 0{,}406), \\
         &        &   & disapproval (L6, 0{,}521), caring (L6, 0{,}502), \\
         &        &   & surprise (L7, 0{,}558), confusion (L7, 0{,}502), \\
         &        &   & realization (L7, 0{,}327) \\
\addlinespace
Tardío   & L8--L11 & 6 & joy (L8, 0{,}601), desire (L8, 0{,}582), \\
         &         &   & approval (L8, 0{,}547), excitement (L8, 0{,}472), \\
         &         &   & annoyance (L8, 0{,}416), \\
         &         &   & disappointment (L9, 0{,}361) \\
\bottomrule
\end{tabular}
\end{table}

\begin{figure}[H]
\centering
\includegraphics[width=0.85\textwidth]{figures/cap5/cap5_probing_heatmap_es.png}
\caption{Heatmap de F1 de los probes (23 emociones $\times$ 13 capas). El
marcador indica la capa de cristalización de cada emoción. Las emociones
se ordenan por profundidad de cristalización creciente.}
\label{fig:probing-heatmap}
\end{figure}

La interpretación de estos grupos conecta la profundidad de procesamiento
con la naturaleza lingüística de la señal emocional. Las emociones del
grupo temprano poseen vocabulario altamente distintivo:
\emph{gratitude} se asocia a marcadores léxicos explícitos (``thanks'',
``thank you''), \emph{amusement} a indicadores como ``lol'' o ``haha'',
\emph{love} a un campo semántico bien delimitado. En estos casos, la
primera capa del encoder (L0) ya dispone de suficiente información
léxica para que un clasificador lineal alcance un rendimiento cercano
al máximo.

En el extremo opuesto, las emociones tardías requieren integración
contextual para desambiguarse de categorías vecinas:
\emph{annoyance} debe distinguirse de \emph{anger},
\emph{approval} de \emph{admiration}, \emph{disappointment} de
\emph{sadness}. Estas distinciones no están disponibles en el
vocabulario de superficie y emergen solo tras múltiples capas de
procesamiento atencional.

Las emociones medias ocupan posiciones intermedias y a menudo combinan
señales léxicas con requisitos contextuales. \emph{Embarrassment} es un
caso ilustrativo: tiene cierta señal léxica (``embarrassed'',
``awkward''), pero su uso pragmático (a veces autoirónico, a veces
genuino) demanda contexto, lo que retrasa su cristalización a L5.


La Figura~\ref{fig:cap5-galaxia-evolution} visualiza esta cristalización
progresiva proyectando las representaciones del token \texttt{[CLS]} en
cuatro capas representativas mediante t-SNE.
 
\begin{figure}[H]
    \centering
    \includegraphics[width=\textwidth]{figures/cap5/cap5_galaxia_evolution_es.png}
    \caption{Proyecciones t-SNE de las representaciones del token
    \texttt{[CLS]} en cuatro capas del modelo. Los puntos pequeños son
    ejemplos individuales del conjunto de test; los rombos grandes
    representan los centroides de cada emoción. Los colores agrupan las
    23 emociones en seis clusters semánticos derivados del análisis
    neuronal (\S\ref{sec:neuron-analysis}). El embedding inicial
    apenas diferencia las emociones (los centroides colapsan en una
    región pequeña). En L4 los puntos se distribuyen, pero los
    centroides aún se solapan. En L8 emergen estructuras parciales:
    las positivas de alta energía se separan, las negativas internas
    comienzan a agruparse. En L11 los seis clusters están claramente
    diferenciados.}
    \label{fig:cap5-galaxia-evolution}
\end{figure}
 
La progresión es consistente con los grupos de cristalización que la
Tabla~\ref{tab:cristalizacion-grupos} identifica numéricamente. En el
embedding ninguna emoción es separable, lo que es esperable: el embedding
codifica el significado del token \texttt{[CLS]} antes de procesar la
frase. En L4 se observa el comportamiento medio del cluster ``Positivas
alta energía'' (gratitude, amusement, love), cuyos centroides ya se
desplazan hacia regiones distintivas, anticipando que cristalizan en
L0--L2. En L8 emergen las distinciones contextuales: el cluster de
``Negativas internas'' se concentra en la región superior, el de
``Negativas reactivas'' en la derecha. En L11 los seis clusters ocupan
regiones diferenciadas del plano, lo que cierra la formación
representacional documentada por probing.

\subsubsection{Ganancia informativa por capa}

La contribución incremental de cada capa se cuantifica como la
diferencia en F1 macro medio entre capas consecutivas
(Figura~\ref{fig:probing-info-gain}). Este análisis revela una
arquitectura de procesamiento en dos fases con un detalle
metodológicamente relevante.

\begin{figure}[H]
\centering
\includegraphics[width=0.8\textwidth]{figures/cap5/cap5_information_gain_es.png}
\caption{Ganancia informativa incremental por capa, medida como el
aumento en F1 macro medio respecto a la capa anterior. La capa L0
absorbe el 61\% de la separabilidad final del modelo.}
\label{fig:probing-info-gain}
\end{figure}

La capa de embedding no contiene información emocional (F1 medio = 0
para todas las emociones, valor exacto). La primera capa del encoder
(L0) produce un salto de $+0{,}349$, capturando aproximadamente el 61\%
de la separabilidad lineal final del modelo
($F1_{L0}/F1_{L11} = 0{,}349/0{,}569 \approx 0{,}61$). Este salto
corresponde a la extracción de señal léxica bruta. Las capas L1--L7
contribuyen con ganancias modestas pero positivas
($+0{,}010$ a $+0{,}048$), refinando gradualmente las representaciones.
La capa L8 produce un segundo pico de $+0{,}033$, asociado a la
desambiguación contextual de las emociones tardías. Las capas finales
(L10--L11) aportan ganancias marginales: L10 incluso muestra una
ganancia ligeramente negativa ($-0{,}002$), seguida de un pequeño
$+0{,}004$ en L11.

Esta estructura de dos fases, con captura léxica masiva en L0,
refinamiento contextual con pico en L8 y consolidación marginal en
L10--L11, anticipa un resultado que será central en el resto del
capítulo: la distinción entre crear información emocional y
formatearla para el clasificador. Las capas L10--L11 apenas
crean información nueva según el probing, pero \S\ref{sec:logit-lens}
y \S\ref{sec:activation-patching} demostrarán que son las únicas
capaces de hacer utilizable la información existente para la
cabeza de clasificación. Esta distinción se cuantifica directamente en
la siguiente sección.

\subsubsection{La frecuencia no predice la profundidad}

Un análisis de correlaciones (Tabla~\ref{tab:probing-correlaciones},
Figura~\ref{fig:probing-correlaciones}) descarta la hipótesis de que la
profundidad de cristalización esté determinada por el volumen de datos
de entrenamiento, hipótesis intuitiva pero empíricamente no soportada.

\begin{table}[ht]
\centering
\caption{Correlaciones de Spearman entre propiedades de las emociones
($n = 23$). Solo la correlación entre profundidad de cristalización y
F1 máximo es significativa.}
\label{tab:probing-correlaciones}
\begin{tabular}{lccc}
\toprule
\textbf{Par de variables} & $\rho$ & \textbf{p-valor} & \textbf{Sig.} \\
\midrule
Capa cristalización vs.\ F1 máximo  & $-0{,}673$ & $0{,}0004$ & *** \\
Frecuencia vs.\ capa cristalización & $-0{,}265$ & $0{,}22$   & ns. \\
Frecuencia vs.\ F1 máximo           & $+0{,}33$\phantom{0}  & $0{,}12$   & ns. \\
\bottomrule
\end{tabular}
\end{table}

\begin{figure}[H]
\centering
\includegraphics[width=\textwidth]{figures/cap5/cap5_crystal_correlations_es.png}
\caption{Izquierda: capa de cristalización vs.\ F1 máximo
($\rho = -0{,}673$, $p < 0{,}001$). Derecha: frecuencia en el dataset
vs.\ capa de cristalización ($\rho = -0{,}265$, $p = 0{,}22$).}
\label{fig:probing-correlaciones}
\end{figure}

La correlación entre profundidad de cristalización y F1 máximo es fuerte
y altamente significativa ($\rho = -0{,}673$, $p = 0{,}0004$): las
emociones que cristalizan antes alcanzan mayor rendimiento. Sin embargo,
la frecuencia en el dataset no predice ni la profundidad de
cristalización ($p = 0{,}22$) ni el rendimiento máximo ($p = 0{,}12$).

Un caso ilustrativo: \emph{disgust}, con 793~ejemplos de entrenamiento,
cristaliza en L2; \emph{annoyance}, con 2\,470~ejemplos (un factor
3$\times$), cristaliza en L8. Más datos no adelantan la cristalización.
La profundidad está determinada por la naturaleza semántica de
la emoción (la complejidad de las distinciones lingüísticas que el
modelo necesita construir), no por el volumen de datos disponible.

Esta observación tiene una implicación práctica relevante para la
construcción de datasets de emociones: añadir más ejemplos de emociones
contextualmente complejas no las hace más fáciles de aprender. La
limitación es estructural, no estadística. Si fuera puramente estadística
(falta de datos), \emph{annoyance} cristalizaría más temprano a medida
que aumenta su frecuencia; lo que observamos es lo contrario.

La potencia estadística es limitada por construcción ($n=23$ emociones).
Para acotar la incertidumbre con honestidad, se computa el intervalo de
confianza al 95\% de cada coeficiente mediante \emph{bootstrap por
emoción} (10\,000 réplicas, semilla 42, resampleo de pares emoción
con reemplazo): la correlación cristalización--F1 máximo
($\rho = -0{,}673$) tiene CI95\,$=[-0{,}87, -0{,}33]$, estrictamente
negativo y por tanto significativo; la correlación
frecuencia--cristalización ($\rho = -0{,}265$) tiene
CI95\,$=[-0{,}65, +0{,}19]$, que cruza cero, confirmando estadísticamente
que la no significancia ($p=0{,}22$) refleja ausencia de evidencia, no
evidencia de ausencia. Con $n=23$, el intervalo es necesariamente
amplio; descartar definitivamente la influencia de la frecuencia
requeriría un dataset con más emociones, vía propuesta en
\S\ref{sec:trabajo-futuro}.


% =============================================================
% 5.2 — LOGIT LENS
% =============================================================

\subsection{\emph{Logit lens}: la brecha entre información disponible e información leíble}
\label{sec:logit-lens}

\subsubsection{Diseño experimental y motivación}

El probing del apartado anterior establece cuándo la información
emocional se vuelve linealmente accesible en las representaciones
intermedias del modelo. Sin embargo, el probe es un clasificador
entrenado ad hoc sobre cada capa: detecta información que está
estructurada linealmente en alguna dirección, pero no
necesariamente en la geometría que utiliza la cabeza de clasificación
final del modelo. Una representación puede ser linealmente separable
según el probe y, simultáneamente, ilegible para el clasificador
entrenado, si la dirección de separación no coincide con la base sobre
la que opera la cabeza de clasificación.

Para distinguir información disponible de información
formateada para el clasificador, se aplica la técnica del logit
lens \cite{nostalgebraist2020logit}: tomar la representación del token
\texttt{[CLS]} de cada capa intermedia y aplicarle directamente la
cabeza de clasificación entrenada del modelo, sin re-entrenamiento. Las
probabilidades resultantes indican la predicción que el modelo emitiría
si tuviera que decidir desde esa capa, utilizando la cabeza ya
optimizada para el \texttt{[CLS]} de L11.

\paragraph{Limitación de calibración.}
Esta técnica tiene una limitación conocida que conviene reconocer
explícitamente: la cabeza de clasificación fue entrenada exclusivamente
sobre representaciones de L11 tras LayerNorm, por lo que aplicarla a
capas intermedias produce predicciones cuyas \emph{magnitudes
absolutas} no son directamente comparables al baseline. Sin embargo,
las comparaciones relativas, entre capas para una misma
emoción o entre emociones para una misma capa, preservan información
sobre la alineación geométrica de la representación con la base del
clasificador. Las cifras absolutas funcionan como diagnóstico de tendencias, no como medida calibrada. La forma cualitativa del patrón es robusta a la
descalibración (se corrobora con probing y patching), pero las
magnitudes específicas no deben interpretarse como probabilidades
absolutas. Aproximaciones más rigurosas como \textit{tuned
lens}~\cite{belrose2023eliciting} entrenan una transformación lineal
por capa para corregir el desajuste de distribución; se discute como
trabajo futuro en el Capítulo~\ref{sec:conclusiones}.

El experimento se realiza sobre 2\,300 ejemplos del conjunto de test
(submuestreo equilibrado entre los tres grupos de cristalización),
agregando los resultados según los grupos del probing
(\S\ref{sec:probing}, Tabla~\ref{tab:cristalizacion-grupos}).

\subsubsection{Patrón agregado: tres regímenes geométricos}

\begin{figure}[H]
\centering
\includegraphics[width=0.85\textwidth]{figures/cap5/cap5_iterative_u_es.png}
\caption{\emph{Logit lens} agregado sobre 2\,300 ejemplos del conjunto de test.
Tres curvas: top-1 sigmoid (probabilidad media de la emoción más
probable), gold sigmoid (probabilidad media de la emoción correcta), y
suma de las 23 sigmoids (masa total de probabilidad). El patrón en U
distingue tres regímenes geométricos sucesivos.}
\label{fig:logit-lens-agregado}
\end{figure}

Promediando sobre los 2\,300 ejemplos, la probabilidad asignada a la
emoción correcta (\emph{gold sigmoid}) sigue un patrón en forma de U a
lo largo del modelo (Figura~\ref{fig:logit-lens-agregado}): parte de un
valor inicial moderado ($\approx 0{,}24$ en L0), cae a casi cero en
capas medias ($\approx 0{,}02$ en L7--L8) y resurge bruscamente en
capas tardías hasta alcanzar 0{,}44 en L11. Este patrón distingue tres
regímenes geométricos cualitativamente distintos:

\paragraph{Régimen léxico (Emb--L2).}
El \texttt{[CLS]} permanece próximo al embedding crudo y activa
difusamente múltiples direcciones del clasificador. La probabilidad
media de la top-1 es alta ($\approx 0{,}6$), pero la masa total de las
23 emociones también lo es ($\approx 5{,}4$), reflejando una activación
distribuida sin discriminación. El clasificador detecta señal léxica
sin filtrar, no decisión emocional.

\paragraph{Régimen de transición (L3--L8).}
Las probabilidades colapsan \emph{para todas las emociones
simultáneamente}, incluida la correcta. Este valle universal indica
que el \texttt{[CLS]} reside en una región del espacio
representacional desalineada con la base del clasificador: la
información que el probing detecta como presente en estas capas
(\S\ref{sec:probing}) existe en una \emph{geometría interna del
modelo}, no proyectable directamente sobre las direcciones de
clasificación. Es la fase de cómputo intermedio.

\paragraph{Régimen de cristalización (L9--L11).}
La probabilidad de la emoción correcta aumenta abruptamente, mientras
que la suma total se mantiene moderada ($\approx 1{,}3$ en L11). El
modelo no solo amplifica la señal correcta: también suprime activamente
las incorrectas. La masa total de las 23 probabilidades es de 1{,}3,
indicando concentración de masa en pocas direcciones. Esta es la fase
de proyección sobre la base del clasificador.

La Figura~\ref{fig:logit-lens-grupos-aggregate} desagrega el mismo
patrón en U según los tres grupos de cristalización del probing
(\S\ref{sec:probing}), mostrando que aunque la \emph{forma} es
universal, la \emph{magnitud} de cada régimen depende del tipo de
emoción. Las emociones tempranas alcanzan mayor gold sigmoid en L11
(consolidación más limpia), mientras que las tardías presentan un valle
más profundo en L4--L8 y una recuperación más modesta.

\begin{figure}[H]
\centering
\includegraphics[width=0.85\textwidth]{figures/cap5/cap5_logit_lens_groups_es.png}
\caption{\emph{Logit lens} desagregado por grupo de cristalización emocional.
Las tres curvas reproducen el patrón en U pero con magnitudes y
tiempos distintos: las emociones tempranas (azul) alcanzan los valores
más altos en L11; las tardías (rojo) muestran el valle más profundo
y una recuperación más modesta. La forma universal del patrón sugiere
que el régimen geométrico es propiedad arquitectónica del modelo, no
de la emoción.}
\label{fig:logit-lens-grupos-aggregate}
\end{figure}

\subsubsection{La brecha probing-\emph{logit lens} revela el trabajo de las
capas tardías}

Superponer el F1 del probe (información disponible) con el gold sigmoid
del \emph{logit lens} (información leíble por el clasificador) sobre los
grupos de cristalización emocional revela un patrón sistemático
(Figura~\ref{fig:logit-lens-grupos}): la brecha entre ambas
curvas cuantifica el trabajo de alineación geométrica que realizan las
capas tardías, y su magnitud depende del tipo de emoción.

\begin{figure}[H]
\centering
\includegraphics[width=\textwidth]{figures/cap5/cap5_lens_vs_probe_es.png}
\caption{Comparación entre probe F1 (información disponible) y gold
sigmoid del \emph{logit lens} (información leíble por el clasificador) por
grupo de cristalización emocional. La región sombreada entre ambas
curvas representa información presente en la representación pero no
proyectable sobre la base del clasificador.
Arriba: emociones tempranas ($n = 9$, 900 frases).
Centro: emociones medias ($n = 8$, 800 frases).
Abajo: emociones tardías ($n = 6$, 600 frases).}
\label{fig:logit-lens-grupos}
\end{figure}

\begin{table}[ht]
\centering
\caption{Brecha entre F1 del probe y gold sigmoid del \emph{logit lens}, por
capa y grupo emocional. La brecha mide la información presente en la
representación pero no leíble por la cabeza de clasificación.}
\label{tab:logit-lens-brechas}
\begin{tabular}{lcccccc}
\toprule
\textbf{Grupo} & \textbf{L0} & \textbf{L4} & \textbf{L7} & \textbf{L9} & \textbf{L11}
& \textbf{Brecha máx.} \\
\midrule
Tempranas & 0{,}37 & 0{,}57 & 0{,}63 & 0{,}62 & 0{,}09 & 0{,}63 (en L7) \\
Medias    & 0{,}02 & 0{,}29 & 0{,}50 & 0{,}52 & 0{,}14 & 0{,}52 (en L9) \\
Tardías   & $-0{,}05$ & 0{,}15 & 0{,}35 & 0{,}44 & 0{,}13 & 0{,}44 (en L9) \\
\bottomrule
\end{tabular}
\end{table}

Este patrón demuestra que las capas tardías realizan dos tipos
cualitativamente distintos de trabajo según el contenido:

\paragraph{Para emociones tempranas,}
la información léxica está disponible desde L0 (probe F1 $\approx 0{,}6$
ya en L0, prácticamente plateau hasta L11), pero el clasificador no
puede leerla hasta L9--L11 (\emph{logit lens} en el suelo durante L4--L8 y
resurgir abrupto en L9--L11). La brecha máxima entre ambas curvas
alcanza 0{,}63 en L7 ---justo antes del bloque tardío--- y se cierra
abruptamente en L9--L11. Es decir, las capas L8--L11 realizan, para
estas emociones, un trabajo de alineación geométrica considerable:
rotar información ya formada en capas anteriores hacia la base del
clasificador, sin crear señal nueva.

\paragraph{Para emociones tardías,}
la curva del probe sube progresivamente entre L0 y L9 (de 0{,}20 a
0{,}48, sin plateau temprano), mostrando que la información se
\emph{construye gradualmente} a lo largo de las capas, no está
disponible desde el principio. La curva del \emph{logit lens} sigue el mismo
perfil con cierto retraso, alcanzando el máximo en L11. La brecha
máxima es de solo 0{,}44 en L9. Para estas emociones, formación
y alineación se realizan en paralelo, no en fases separadas.

\paragraph{Las emociones medias}
ocupan una posición intermedia: la información se forma entre L0 y L7
(probe sube de 0{,}32 a 0{,}55), y la alineación se concentra en
L9--L11. Brecha máxima de 0{,}52 en L9.

\subsubsection{Conexión con \emph{activation patching}}

Esta interpretación geométrica explica un resultado central del
\emph{activation patching} que se presenta a continuación
(\S\ref{sec:activation-patching}): por qué la restauración exclusiva de
la capa~11 desde un modelo colapsado recupera el 100\% del rendimiento.
La capa~11 no contiene ``todo el conocimiento emocional'' en sentido
literal (el probing muestra que la información se distribuye a lo
largo de todas las capas), pero es donde culmina la proyección
de esa información sobre la base del clasificador. Restaurar L11
reconstituye la última etapa de alineación geométrica, y la información
residual que sobrevive a las capas comprimidas anteriores resulta
suficiente para que el clasificador produzca predicciones correctas.

Esto reconcilia tres hallazgos aparentemente disonantes que se
desarrollarán en las secciones siguientes:

\begin{enumerate}[nosep]
\item el probing localiza la formación de información en capas
tempranas y medias;
\item el logit lens muestra que esa información no es leíble hasta
capas tardías;
\item el activation patching identifica L11 como suficiente para la
recuperación.
\end{enumerate}

Las tres mediciones son consistentes si se distingue entre tres
conceptos relacionados pero distintos: representación (qué
información está presente, medido por probing), alineación
(esa información está proyectada sobre la base del clasificador, medido
por \emph{logit lens}) y acceso (qué componentes son suficientes
para que el clasificador opere, medido por patching). Esta clarificación
tripartita es, en sí misma, una contribución metodológica del trabajo:
sugiere que cualquier análisis interpretativo que confunda los tres
conceptos producirá conclusiones contradictorias como las que el
probing y el patching parecen producir aisladamente.
% =============================================================
% 5.3 — ACTIVATION PATCHING
% =============================================================

\subsection{\emph{Activation patching}: localización del conocimiento emocional}
\label{sec:activation-patching}

\subsubsection{Diseño experimental}

El probing identifica dónde se forma la información emocional;
el \emph{logit lens} identifica dónde se alinea con la base del
clasificador; el \emph{activation patching} identifica \emph{qué componentes
son funcionalmente necesarios} para que el modelo opere. Las tres
preguntas son complementarias: una capa puede crear poca información
nueva (bajo gain en probing) pero ser imprescindible para que la
información existente llegue al clasificador en formato utilizable.

El experimento parte del modelo comprimido con SVD uniforme a rango 64,
cuyo rendimiento es F1 macro = 0{,}000 (colapso total documentado en
\S\ref{sec:compresion-uniforme}). Desde este estado de destrucción
completa, se restauran los pesos originales de una capa o componente
a la vez y se mide la recuperación del F1 macro. El porcentaje de
restauración se calcula como
$\text{F1}_{\text{patched}} / \text{F1}_{\text{baseline}} \times 100$,
donde $\text{F1}_{\text{baseline}} = 0{,}577$.

\paragraph{Naturaleza funcional vs.\ causal.}
El sentido en el que este experimento establece causalidad merece una matización: la corrupción mediante SVD truncada no es neutral
en sentido estadístico (a diferencia del \emph{causal tracing} clásico de
Meng et al.~\cite{meng2022locating}, que usa noise gaussiano), sino una
\emph{distorsión estructurada}. Esto significa que las conclusiones
que extraemos son de naturaleza funcional (qué componentes
son suficientes para reactivar el modelo desde el colapso) más que
estrictamente causal en el sentido de Pearl. La distinción es
relevante para la replicabilidad: el \emph{activation patching} aquí responde
a la pregunta ``¿qué componente del modelo original es suficiente
para superar la perturbación SVD?'', no a la pregunta ``¿qué
componente codifica causalmente el conocimiento emocional?''. La
distinción no debilita los resultados ---los hace más precisos en su
alcance epistémico.

\subsubsection{Patching por capa completa}

La restauración capa por capa (Tabla~\ref{tab:patching-capas}) revela
una concentración extrema de la capacidad funcional en las capas
finales.

\begin{table}[ht]
\centering
\caption{Restauración del F1 macro al parchear una sola capa del modelo
comprimido (rango 64). Se omiten las capas con restauración nula
(0--7).}
\label{tab:patching-capas}
\begin{tabular}{ccc}
\toprule
\textbf{Capa restaurada} & \textbf{F1 Macro} & \textbf{Restauración (\%)} \\
\midrule
0--7 & 0{,}000 & 0{,}0 \\
8    & 0{,}001 & 0{,}1 \\
9    & 0{,}024 & 4{,}1 \\
10   & 0{,}139 & 24{,}1 \\
11   & 0{,}577 & 100{,}0 \\
\bottomrule
\end{tabular}
\end{table}

Las capas 0--7 no contribuyen en absoluto a la recuperación. La capa~11
sola es \emph{suficiente} para restaurar el 100\% del rendimiento
baseline. La interpretación correcta de este resultado, anticipada por
el \emph{logit lens} (\S\ref{sec:logit-lens}), es que L11 es donde culmina la
alineación geométrica de la representación con la base del
clasificador: restaurar L11 reconstituye la última proyección, y la
señal residual que sobrevive a 10 capas comprimidas resulta suficiente
para que el clasificador produzca predicciones correctas.

Una posible objeción al resultado de patching de L11 conviene anticiparla.
Podría argumentarse que el clasificador está conectado precisamente a
L11 mediante una capa lineal entrenada, y que por tanto la
``suficiencia'' de L11 es trivial. La objeción no aplica, por dos motivos. Primero, el clasificador no está conectado solo a la salida
cruda de L11 sino a L11 \emph{tras LayerNorm} y, en BERT, tras la
cabeza \textit{pooler} (CLS de L11 $\rightarrow$ tanh $\rightarrow$
linear). Restaurar los pesos de L11 incluye las seis matrices del
encoder pero \emph{no} la cabeza pooler ni el clasificador, que ya
estaban no comprimidos en el modelo \textit{corrupted run}. El hallazgo
no es ``el clasificador ve L11'' (eso sería trivial) sino ``L11 tras 10 capas comprimidas todavía produce señal suficiente para que el
clasificador clasifique correctamente'', lo cual no es trivial.
Segundo, el \emph{activation patching} sobre las capas 8--10 muestra
restauración parcial pero no completa, lo que indica que la suficiencia
de L11 no es un artefacto de localización: cuando otras capas
``cercanas'' al clasificador se restauran, la recuperación es parcial.
La especificidad funcional de L11 es genuina.

\subsubsection{La FFN como locus computacional crítico}

El análisis por componente (Tabla~\ref{tab:patching-componentes})
desglosa cada capa en sus dos subcomponentes (atención multi-cabeza y
red feed-forward) y revela una asimetría funcional de primer
orden.

\begin{table}[ht]
\centering
\caption{Restauración del F1 macro al parchear un solo componente
(atención o FFN) de las capas 8--11 del modelo comprimido.}
\label{tab:patching-componentes}
\begin{tabular}{cccc}
\toprule
\textbf{Capa} & \textbf{Componente} & \textbf{F1 Macro} & \textbf{Restauración (\%)} \\
\midrule
11 & FFN       & 0{,}577 & \textbf{100{,}0} \\
11 & Atención  & 0{,}365 & 63{,}3 \\
10 & Atención  & 0{,}156 & 27{,}1 \\
10 & FFN       & 0{,}139 & 24{,}1 \\
9  & Atención  & 0{,}033 & 5{,}6 \\
9  & FFN       & 0{,}024 & 4{,}1 \\
8  & Atención  & 0{,}000 & 0{,}0 \\
8  & FFN       & 0{,}001 & 0{,}1 \\
\bottomrule
\end{tabular}
\end{table}

El resultado más destacado es que restaurar únicamente la FFN
de la capa~11, aproximadamente un tercio de los parámetros de esa
capa, recupera el 100\% del rendimiento. La atención de la misma
capa solo alcanza un 63{,}3\%, y con alta variabilidad entre emociones.
Este hallazgo posiciona a la FFN de la capa~11 como el locus
computacional crítico de la capacidad emocional del modelo.

La dominancia de la FFN desafía productivamente la narrativa centrada
en la atención que ha prevalecido desde \emph{Attention Is All You
Need} \cite{vaswani2017attention}. En este modelo \emph{fine-tuned} para
clasificación emocional, la atención contribuye a enrutar información
entre posiciones, pero es la FFN la que realiza la transformación
computacional que convierte representaciones contextualizadas en
representaciones clasificables. Este resultado es coherente con
trabajos recientes que identifican las FFN como \textit{memorias
asociativas} del modelo \cite{geva2021transformer}, y conecta
directamente con el hallazgo del \emph{logit lens}: la proyección sobre la
base del clasificador que documentamos en \S\ref{sec:logit-lens} es
realizada principalmente por las FFN de las capas tardías, no por la
atención.

\subsubsection{Restauración diferencial por emoción}

La restauración de la atención de la capa~11 presenta una marcada
heterogeneidad entre emociones (Tabla~\ref{tab:patching-emociones}).

\begin{table}[ht]
\centering
\caption{Restauración del F1 por emoción al parchear solo la atención
de la capa~11. Se muestran las 5 emociones mejor y peor restauradas.}
\label{tab:patching-emociones}
\begin{tabular}{lclc}
\toprule
\textbf{Bien restauradas} & \textbf{Rest.\ (\%)} &
\textbf{Mal restauradas} & \textbf{Rest.\ (\%)} \\
\midrule
gratitude   & 100{,}0 & remorse        & 0{,}0 \\
annoyance   & 98{,}8  & fear           & 3{,}6 \\
admiration  & 95{,}4  & desire         & 4{,}4 \\
anger       & 94{,}5  & realization    & 10{,}0 \\
curiosity   & 91{,}1  & disappointment & 18{,}8 \\
\bottomrule
\end{tabular}
\end{table}

Las emociones que la atención restaura bien son aquellas con señales
léxicas claras: \emph{gratitude} se identifica por marcadores
explícitos, \emph{anger} por vocabulario de alta valencia negativa. La
atención puede redirigir estas señales al token \texttt{[CLS]} de
forma eficaz incluso con FFN comprimidas. En cambio, las emociones
que la atención no restaura (\emph{remorse}, \emph{fear},
\emph{desire}) son aquellas que requieren computación contextual
más rica, precisamente la función que desempeña la FFN.

Este patrón refuerza la interpretación de los grupos de
cristalización del probing (\S\ref{sec:probing}) y del \emph{logit lens}
(\S\ref{sec:logit-lens}): las emociones tempranas son léxicas y
dependen menos de la FFN profunda; las emociones tardías son
contextuales y necesitan la FFN intacta para completar tanto la
formación como la alineación de su representación.


% =============================================================
% 5.4 — ABLACIÓN DE CABEZAS DE ATENCIÓN
% =============================================================

\subsection{Ablación de cabezas de atención}
\label{sec:head-ablation}

\subsubsection{Diseño experimental}

Los experimentos anteriores operan a nivel de capa completa o de
subcomponente (atención vs.\ FFN). La ablación de cabezas desciende a
la granularidad individual: para cada una de las 144~cabezas de
atención (12~capas $\times$ 12~cabezas), se anulan sus pesos
(\textit{zero-out}) y se mide la caída en F1 macro y en F1 por
emoción. Para mantener tractable el coste computacional ---144~cabezas
$\times$ una pasada de evaluación cada una--- la ablación se ejecuta
sobre un submuestreo de 2\,000 ejemplos del conjunto de test (en lugar
del test completo de 5\,427). El baseline local sobre este submuestreo
es 0{,}582, ligeramente superior al baseline global de 0{,}577 del
test completo (\S\ref{sec:met-modelo}); todas las caídas reportadas en
esta sección son absolutas respecto a ese baseline local, lo que las
mantiene comparables entre sí.

A partir de estas mediciones se construyen dos métricas ortogonales:
la \emph{importancia}, definida como el impacto agregado en F1 a
través de todas las emociones, y la \emph{selectividad}, que mide la
concentración de dicho impacto en pocas emociones (valor alto = efecto
concentrado en 1--2 emociones; valor bajo = efecto distribuido). La
ortogonalidad de ambas dimensiones se verifica empíricamente
($r = -0{,}055$), lo que habilita una clasificación 2$\times$2 de las
cabezas. La consistencia interna de la métrica de importancia se
valida mediante su correlación con la caída real de F1 por ablación
($r = 0{,}71$).

\subsubsection{Taxonomía funcional}

Las 144~cabezas se clasifican en cuatro categorías
(Tabla~\ref{tab:head-taxonomy}).

\begin{table}[ht]
\centering
\caption{Taxonomía de cabezas de atención basada en importancia y
selectividad. La baja correlación entre ambas dimensiones
($r = -0{,}055$) valida la clasificación.}
\label{tab:head-taxonomy}
\begin{tabular}{lcccl}
\toprule
\textbf{Categoría} & \textbf{N} & \textbf{Import.} & \textbf{Select.} & \textbf{Rol funcional} \\
\midrule
Critical Specialist & 37 & 0{,}282 & 0{,}680 & Especializada en emociones concretas \\
Critical Generalist & 35 & 0{,}265 & 0{,}544 & Infraestructura emocional general \\
Minor Specialist    & 35 & 0{,}101 & 0{,}702 & Nicho funcional (1--2 emociones) \\
Dispensable         & 37 & 0{,}120 & 0{,}556 & Candidata a poda \\
\bottomrule
\end{tabular}
\end{table}

\subsubsection{Gradiente de criticidad por profundidad}

El hallazgo principal de este análisis es que la proporción de cabezas
críticas crece monotónicamente con la profundidad
(Tabla~\ref{tab:head-profundidad}, Figura~\ref{fig:head-taxonomy-depth}).

\begin{table}[ht]
\centering
\caption{Distribución de cabezas críticas y dispensables por bloque de
profundidad. La capa 11 no contiene ninguna cabeza dispensable.}
\label{tab:head-profundidad}
\begin{tabular}{lcc}
\toprule
\textbf{Bloque} & \textbf{Críticas} & \textbf{Dispensables} \\
\midrule
Capas 0--4  & 16/60 (27\%) & 22/60 (37\%) \\
Capas 5--7  & 19/36 (53\%) & 10/36 (28\%) \\
Capas 8--11 & 37/48 (77\%) & \phantom{0}5/48 (10\%) \\
\bottomrule
\end{tabular}
\end{table}

La capa~11 constituye un caso extremo: ninguna de sus 12~cabezas es
dispensable: las 12 son críticas (6~specialists y 6~generalists). La
importancia media por capa presenta un salto abrupto a partir de la
capa~8 ($\sim 0{,}25$ vs.\ $\sim 0{,}15$ en capas anteriores), con la
capa~11 alcanzando una importancia media de 0{,}381. Las 5~cabezas más
importantes del modelo pertenecen a las capas 10--11, encabezadas por
L11-H11 con una caída de F1 macro de 0{,}023.

\begin{figure}[H]
\centering
\includegraphics[width=0.85\textwidth]{figures/cap5/cap5_head_categories_es.png}
\caption{Distribución de categorías de cabezas por capa. La proporción
de cabezas críticas (specialists + generalists) crece con la
profundidad. La capa~11 no contiene ninguna cabeza dispensable.}
\label{fig:head-taxonomy-depth}
\end{figure}

Este resultado proporciona triangulación con cuatro métodos
independientes: el probing (\S\ref{sec:probing}) muestra que la
separabilidad lineal máxima se alcanza en capas finales; el \emph{logit lens}
(\S\ref{sec:logit-lens}) identifica L9--L11 como la fase de
cristalización donde la representación se alinea con el clasificador;
el \emph{activation patching} (\S\ref{sec:activation-patching}) demuestra
que la capa~11 es suficiente para restaurar el rendimiento completo;
y la ablación confirma que cada cabeza individual de esa capa es
funcionalmente necesaria. Cinco técnicas distintas convergen en la
misma conclusión: las capas 8--11 son el núcleo funcional del modelo
emocional.

\subsubsection{Cabezas interferentes}

Un hallazgo secundario es la existencia de 21~cabezas cuya ablación
\emph{mejora} el rendimiento del modelo (caída de F1 más negativa que
$-0{,}003$). Aplicando un umbral más laxo ($\Delta F1 < 0$) el conteo
sube a 79 cabezas, pero la mayoría de esos efectos están en el rango
del ruido de medición; las 21~cabezas que superan el umbral de
$0{,}003$ son las únicas con un beneficio cuantificable. La
Tabla~\ref{tab:cabezas-interferentes} lista las diez con mayor efecto.

\begin{table}[ht]
\centering
\caption{Top 10 cabezas interferentes (ordenadas por mejora del F1
macro tras ablación). Categoría según la taxonomía de
\S\ref{sec:head-ablation}.}
\label{tab:cabezas-interferentes}
\small
\begin{tabular}{c c l c c}
\toprule
\textbf{Capa} & \textbf{Cabeza} & \textbf{Categoría} & \textbf{$\Delta$F1 (mejora)} & \textbf{Selectividad} \\
\midrule
11 &  5 & Critical Specialist  & $+$0{,}0121 & 0{,}635 \\
10 & 11 & Critical Generalist  & $+$0{,}0077 & 0{,}596 \\
 1 & 10 & Critical Specialist  & $+$0{,}0051 & 0{,}674 \\
 2 &  1 & Critical Generalist  & $+$0{,}0047 & 0{,}469 \\
 3 &  7 & Critical Specialist  & $+$0{,}0047 & 0{,}750 \\
 7 &  7 & Minor Specialist     & $+$0{,}0046 & 0{,}829 \\
 7 &  1 & Critical Generalist  & $+$0{,}0045 & 0{,}583 \\
 5 &  1 & Critical Specialist  & $+$0{,}0043 & 0{,}673 \\
 7 &  9 & Critical Specialist  & $+$0{,}0043 & 0{,}703 \\
 5 &  2 & Critical Specialist  & $+$0{,}0043 & 0{,}621 \\
\bottomrule
\end{tabular}
\end{table}

La distribución por profundidad es informativa: 7 interferentes en
capas tempranas (L0--L4), 9 en intermedias (L5--L7) y solo 5 en
tardías (L8--L11). El 76\% se concentra en capas no críticas para la
clasificación final, lo que es coherente con la taxonomía
2$\times$2: las cabezas dispensable y las críticas-generalistas son
las que más contribuyen al ruido cuando se acumulan. La cabeza más
interferente es L11-H5 ($\Delta F1 = +0{,}012$), que además es
nominalmente \emph{Critical Specialist}; este caso aparentemente
paradójico sugiere que su contribución específica entra en conflicto con
otras cabezas de L11, y eliminarla resuelve la interferencia.

Como estimación aditiva, ablacionar simultáneamente las 21~cabezas
interferentes daría una ganancia teórica de $+0{,}095$ en F1 macro
(suma directa de las caídas negativas), lo que llevaría el modelo de
0{,}577 a $\sim$0{,}672. El supuesto de aditividad es probablemente
demasiado optimista debido a interacciones no-lineales entre cabezas
ablacionadas, pero incluso una fracción significativa de esa ganancia
sería notable para una intervención sin re-entrenamiento. El
experimento de ablación coordinada se propone como trabajo futuro
(\S\ref{sec:trabajo-futuro}).

Este hallazgo es, en sí mismo, una observación metodológicamente
relevante: incluso un modelo \emph{fine-tuned} eficaz contiene componentes
contraproducentes, lo que abre la puerta a optimizaciones post-hoc
que no requieren acceso al proceso de entrenamiento ---ni siquiera a
los gradientes--- y se complementan naturalmente con la compresión
SVD del Capítulo~\ref{sec:resultados-compresion}.

\subsubsection{Especialización emocional}

Cada emoción depende de cabezas diferentes, concentradas
mayoritariamente en las capas 8--11
(Tabla~\ref{tab:head-emociones}).

\begin{table}[ht]
\centering
\caption{Cabeza más importante para cada emoción (top 10 por magnitud
de impacto). La caída de F1 indica cuánto pierde esa emoción al
ablacionar su cabeza más crítica.}
\label{tab:head-emociones}
\begin{tabular}{lcc}
\toprule
\textbf{Emoción} & \textbf{Cabeza crítica} & \textbf{$\Delta$F1} \\
\midrule
embarrassment  & L10-H1 & $-0{,}320$ \\
annoyance      & L10-H0 & $-0{,}090$ \\
disappointment & L10-H4 & $-0{,}080$ \\
realization    & L11-H6 & $-0{,}064$ \\
fear           & L0-H8  & $-0{,}056$ \\
confusion      & L9-H2  & $-0{,}044$ \\
sadness        & L11-H6 & $-0{,}041$ \\
disgust        & L9-H7  & $-0{,}040$ \\
joy            & L10-H7 & $-0{,}039$ \\
anger          & L8-H8  & $-0{,}031$ \\
\bottomrule
\end{tabular}
\end{table}

El caso de \emph{embarrassment} es extremo: pierde un 32\% de su F1
al ablacionar cualquiera de cuatro cabezas equivalentes
(L10-H1, L10-H4, L11-H2 y L11-H8, todas con $\Delta\text{F1} = -0{,}320$
por empate exacto en el ranking). Esta concentración funcional en un
puñado pequeño de cabezas redundantes ---no en una sola, pero sí en
muy pocas--- probablemente refleja la baja frecuencia de la emoción
en GoEmotions (303~ejemplos), que hace que el modelo aprenda un
circuito atencional acotado y replicado por seguridad estadística.
Desde la perspectiva de la compresión, esto implica que la eliminación
o degradación coordinada de esas cuatro cabezas destruiría
selectivamente la capacidad del modelo para detectar
\emph{embarrassment}, un ejemplo de sesgo silencioso por compresión
que las métricas agregadas no detectan fácilmente.

Otro patrón relevante es el compartimiento de circuitos:
\emph{realization} y \emph{sadness} dependen de la misma cabeza
(L11-H6), lo que sugiere un mecanismo atencional común
, posiblemente la detección de un patrón contextual de expectativa
frustrada que ambas emociones comparten.


% =============================================================
% 5.5 — ANÁLISIS NEURONAL DE LAS FFN
% =============================================================

\subsection{Análisis neuronal de las redes feed-forward}
\label{sec:neuron-analysis}

\subsubsection{Diseño experimental}

El \emph{activation patching} estableció la FFN de la capa~11 como el locus
computacional crítico. El análisis neuronal desciende al nivel de las
neuronas individuales dentro de las FFN para examinar dos preguntas:
cómo se distribuye la especialización emocional a través de las
36\,864~neuronas (12~capas $\times$ 3\,072~neuronas), y si la
compresión SVD daña selectivamente las neuronas más especializadas.

La selectividad de cada neurona para cada emoción se cuantifica
mediante un estadístico tipo Cohen's $d$: la diferencia entre la
activación media cuando la emoción está presente y cuando está
ausente, normalizada por la desviación típica agrupada. Se extraen
las activaciones post-GELU de la capa intermedia de cada FFN
(\texttt{BertIntermediate}) mediante \textit{forward hooks}, con
mean-pooling sobre posiciones de secuencia. Se definen como
\emph{significativas} las neuronas con $|d| > 2{,}0$.

Para investigar la conexión con la compresión, se proyecta el vector
de selectividad de cada emoción (un vector de 3\,072~dimensiones por
capa) sobre los subespacios SVD conservado y eliminado de la matriz
$W_{\text{int}}$ a rango 128, y se calcula la fracción de señal
emocional que cae en el subespacio eliminado.

\subsubsection{Especialización concentrada en capas profundas}

La distribución de neuronas significativas por profundidad
(Tabla~\ref{tab:neuron-profundidad}) replica el patrón ya observado
en probing, patching y ablación de cabezas.

\begin{table}[ht]
\centering
\caption{Distribución de neuronas especializadas ($|d| > 2{,}0$) por
bloque de profundidad. El 84\% de las neuronas significativas se
concentra en las capas 8--11.}
\label{tab:neuron-profundidad}
\begin{tabular}{lcc}
\toprule
\textbf{Bloque} & \textbf{Neuronas significativas} & \textbf{Selectividad máx.\ media} \\
\midrule
Capas 0--3  & 11 (0{,}3\%)    & 0{,}91--1{,}15 \\
Capas 4--7  & 570 (16\%)      & 1{,}50--2{,}19 \\
Capas 8--11 & 3\,061 (84\%)   & 3{,}10--3{,}92 \\
\bottomrule
\end{tabular}
\end{table}

El 84\% de las neuronas especializadas se concentra en las capas
8--11, con la capa~11 aportando 1\,127~neuronas significativas por sí
sola. Las capas 0--3 contienen prácticamente ninguna neurona con
especialización emocional detectable. Este resultado es coherente con
la función de las capas tempranas como procesadores léxicos genéricos:
contribuyen a la señal emocional (probing), pero lo hacen a través
de \emph{representaciones distribuidas}, no de neuronas individuales
especializadas.

\subsubsection{Asimetría radical entre emociones}

El número de neuronas significativas varía en órdenes de magnitud
entre emociones (Tabla~\ref{tab:neuron-emociones}).

\begin{table}[ht]
\centering
\caption{Neuronas significativas y selectividad máxima por emoción
(selección). Tres emociones no tienen ninguna neurona con
$|d| > 2{,}0$.}
\label{tab:neuron-emociones}
\begin{tabular}{lcc}
\toprule
\textbf{Emoción} & \textbf{Neuronas sig.} & \textbf{Selectividad máx.} \\
\midrule
gratitude      & 818 & 6{,}88 \\
remorse        & 442 & 7{,}75 \\
love           & 399 & 5{,}41 \\
amusement      & 394 & 5{,}90 \\
fear           & 306 & 7{,}71 \\
\midrule
annoyance      & 0   & 1{,}71 \\
disappointment & 0   & 1{,}89 \\
realization    & 0   & 1{,}96 \\
\bottomrule
\end{tabular}
\end{table}

Las emociones mejor codificadas (\emph{gratitude}, \emph{remorse})
disponen de cientos de neuronas dedicadas, mientras que
\emph{annoyance}, \emph{disappointment} y \emph{realization} carecen
de neuronas individuales significativas. Las neuronas con mayor
selectividad de todo el modelo ($d > 7$) pertenecen a \emph{fear} y
\emph{remorse}, ambas en la capa~11.

Esta asimetría sugiere que el modelo emplea estrategias de
codificación cualitativamente diferentes: las emociones ``fuertemente
escritas'' utilizan neuronas altamente especializadas, mientras que
las emociones ``débilmente escritas'' se codifican mediante patrones
distribuidos más sutiles que la métrica de selectividad individual no
captura.

\subsubsection{La hipótesis SVD-neuronas: un resultado nulo informativo}

La hipótesis principal de este análisis era que la compresión SVD
dañaría selectivamente las emociones cuyas neuronas especializadas
cayeran desproporcionadamente en el subespacio eliminado. Esta
hipótesis resulta \emph{falsa}: la correlación entre la fracción de
señal en el subespacio eliminado y la caída de F1 por emoción es
prácticamente nula (Spearman $\rho = -0{,}004$, $p = 0{,}99$).

Además, la fracción de señal emocional que cae en el subespacio
eliminado es prácticamente idéntica para todas las emociones (rango
0{,}53--0{,}56, media 0{,}54). La SVD no discrimina entre señales
emocionales: elimina señal de forma uniforme, independientemente de
la identidad funcional de las neuronas afectadas.

Este resultado nulo es informativo porque revela una disociación clara: la geometría que la SVD optimiza (energía de valores
singulares) es ortogonal a la estructura funcional del modelo
(especialización emocional). La compresión opera en un espacio
matemático que no ``ve'' la función. Esto es teóricamente importante
porque clarifica por qué la compresión SVD informada por
interpretabilidad funciona: no porque la SVD pueda dirigirse a
neuronas específicas, sino porque puede asignar más rango a los
componentes (FFN tardía) cuya destrucción causa más daño.

Este patrón es consistente con la hipótesis de \emph{superposition}
formalizada por Elhage et al.~\cite{elhage2022toy}: cuando una red
necesita representar más \emph{features} funcionales que dimensiones
disponibles, las distribuye en direcciones casi-ortogonales del
espacio de activación, aceptando pequeñas interferencias a cambio de
mayor capacidad. En esta lectura, las features emocionales en BERT
estarían codificadas en direcciones que no coinciden con las
direcciones SVD principales de los pesos: las primeras provienen del
optimizador de la tarea durante el \emph{fine-tuning}; las segundas
provienen de la geometría algebraica de los pesos. No hay razón a
priori para que coincidan, y nuestro resultado empírico (fracción de
señal eliminada uniforme entre emociones) es exactamente lo que
predice superposition. La implicación práctica es que los métodos de
compresión geométrica como la SVD truncada nunca podrán ser
\emph{selectivamente} funcionales por sí solos: la selectividad
funcional debe entrar por la puerta de la sensibilidad empírica
medida ---es decir, justo lo que el algoritmo greedy del
Capítulo~\ref{sec:sintesis} hace al usar las caídas de F1 por
componente como costes de optimización.

\subsubsection{La norma de selectividad como predictor de vulnerabilidad}

Si la geometría SVD no explica la vulnerabilidad diferencial, ¿qué la
explica? Un análisis sistemático testa seis predictores posibles de
la caída de F1 bajo compresión
(Tabla~\ref{tab:neuron-predictores}).

\begin{table}[ht]
\centering
\caption{Predictores de la vulnerabilidad emocional a compresión SVD
($n = 23$). La norma de selectividad, que mide la intensidad total
de codificación neuronal, es el mejor predictor.}
\label{tab:neuron-predictores}
\begin{tabular}{lccc}
\toprule
\textbf{Predictor} & $\rho$ & \textbf{p-valor} & \textbf{Sig.} \\
\midrule
Fracción en subespacio eliminado & $-0{,}004$ & $0{,}986$ & ns. \\
Nº ejemplos de entrenamiento     & $+0{,}174$ & $0{,}427$ & ns. \\
Aislamiento representacional     & $+0{,}452$ & $0{,}031$ & * \\
Selectividad máxima              & $+0{,}517$ & $0{,}012$ & * \\
Nº neuronas significativas       & $+0{,}581$ & $0{,}004$ & ** \\
\textbf{Norma de selectividad}   & $\mathbf{+0{,}639}$ & $\mathbf{0{,}001}$ & *** \\
\bottomrule
\end{tabular}
\end{table}

El mejor predictor es la \emph{norma del vector de selectividad}
($\rho = 0{,}639$, $p = 0{,}001$), que mide la intensidad total con la
que el modelo codifica cada emoción a través de sus neuronas FFN. Las
emociones con codificación más intensa (\emph{gratitude}: norma 140,
\emph{remorse}: 111, \emph{love}: 105) sufren más con la compresión
que las de codificación débil (\emph{disappointment}: 42,
\emph{realization}: 43).

La interpretación es la siguiente: la SVD no ataca selectivamente las
neuronas emocionales: la geometría del subespacio eliminado es
ortogonal a la especialización funcional, como se ha demostrado. Pero
las emociones que el modelo codifica con mayor intensidad neuronal
son más vulnerables a perturbaciones generales de los pesos,
porque su representación depende de configuraciones más específicas y
exigentes. Es una vulnerabilidad por intensidad, no por
geometría. La distinción es útil metodológicamente: sugiere
que la robustez a perturbación general (no solo a SVD) puede
estimarse a partir de la intensidad de codificación, una métrica
intrínseca al modelo que no requiere experimentos de compresión.

\subsubsection{Clusters emocionales emergentes}

La agrupación jerárquica sobre los vectores de selectividad
concatenados produce seis clusters emocionales
(Tabla~\ref{tab:neuron-clusters}) que reflejan categorías psicológicas
coherentes.

\begin{table}[ht]
\centering
\caption{Clusters emocionales emergentes de la estructura neuronal del
modelo. Las agrupaciones se obtienen por clustering jerárquico sobre
los vectores de selectividad de 36\,864~neuronas.}
\label{tab:neuron-clusters}
\begin{tabular}{cl}
\toprule
\textbf{Cluster} & \textbf{Emociones} \\
\midrule
Positivas de alta energía & admiration, amusement, excitement, gratitude, joy, love \\
Negativas reactivas       & anger, annoyance, disappointment, disapproval, disgust \\
Negativas internas        & embarrassment, fear, remorse, sadness \\
Epistémicas               & confusion, curiosity, surprise \\
Orientadas al otro        & caring, desire, optimism \\
Baja especificidad        & approval, realization \\
\bottomrule
\end{tabular}
\end{table}

Estos clusters no se imponen \emph{a priori} sino que emergen
directamente de la estructura de activaciones del modelo. Que
reproduzcan categorías reconocibles de la psicología de las
emociones (positivas vs.\ negativas, reactivas vs.\
internas, epistémicas vs.\ afectivas) sugiere que el \emph{fine-tuning} ha
organizado las representaciones internas de BERT de forma
psicológicamente coherente, no solo estadísticamente eficaz. Esta
correspondencia es no trivial: el modelo no ha sido entrenado con
información sobre la estructura psicológica de las emociones, y sin
embargo la recupera implícitamente.


% =============================================================
% 5.6 — SÍNTESIS DEL CAPÍTULO
% =============================================================

\subsection{Síntesis: una arquitectura funcional convergente}
\label{sec:interp-sintesis}

Las cinco técnicas de este capítulo emplean métodos diferentes
(clasificadores lineales entrenados \textit{ad hoc}, aplicación
directa del clasificador final, restauración selectiva de pesos,
ablación individual de cabezas, estadísticas de activación neuronal)
pero convergen en una descripción unificada de cómo BERT procesa la
información emocional. La Tabla~\ref{tab:convergencia-tecnicas}
resume esta convergencia.

\begin{table}[ht]
\centering
\caption{Convergencia de cinco técnicas de interpretabilidad en las
mismas conclusiones estructurales.}
\label{tab:convergencia-tecnicas}
\small
\begin{tabular}{p{2.8cm}p{3.5cm}p{6.5cm}}
\toprule
\textbf{Técnica} & \textbf{Pregunta} & \textbf{Resultado clave} \\
\midrule
Probing &
¿Dónde se \emph{forma} la señal? &
L0 captura el 61\% (léxica); L8 segundo pico (contextual).
Capas finales aportan $< 1\%$. \\
\addlinespace
\emph{Logit lens} &
¿Dónde es leíble por el clasificador? &
Patrón en U: alta en L0--L2, valle en L3--L8, cristalización
abrupta en L9--L11. La brecha probing--\emph{logit lens} cuantifica el
trabajo de alineación. \\
\addlinespace
\emph{Activation patching} &
¿Dónde \emph{reside} el conocimiento? &
Solo L11 restaura el 100\%. La FFN sola basta;
la atención solo recupera el 63\%. \\
\addlinespace
Ablación de cabezas &
¿Qué cabezas son \emph{necesarias}? &
77\% críticas en capas 8--11; 0 dispensables en capa~11.
Gradiente monotónico con profundidad. \\
\addlinespace
Análisis neuronal &
¿Qué neuronas están \emph{especializadas}? &
84\% de neuronas significativas en capas 8--11.
La norma de selectividad ($\rho = 0{,}64$) predice vulnerabilidad. \\
\bottomrule
\end{tabular}
\end{table}

De esta convergencia se extraen cuatro conclusiones principales que
constituyen la base sobre la que el Capítulo~\ref{sec:sintesis}
construirá la estrategia de compresión informada.

\paragraph{Arquitectura de tres niveles.}
BERT procesa la información emocional en tres fases cualitativamente
distintas. Las capas tempranas (Emb--L2) extraen señal léxica bruta:
información linealmente separable pero distribuida sobre muchas
direcciones simultáneamente. Las capas intermedias (L3--L8) realizan
cómputo de transición en una geometría \emph{interna} del modelo: la
información se refina y reorganiza, pero permanece desalineada con la
base del clasificador (valle del \emph{logit lens}). Las capas profundas
(L9--L11) ejecutan la cristalización: alinean la representación con
la cabeza de clasificación y suprimen activamente las opciones
incorrectas. Este último paso es el cuello de botella geométrico del modelo, y su criticidad bajo compresión SVD se
explica por su baja redundancia espectral en las matrices que
realizan la alineación.

\paragraph{La paradoja información-dependencia, resuelta.}
El probing muestra que las capas finales (L10--L11) apenas crean
información nueva (gain $< 0{,}004$ del total), pero el activation
patching y la ablación demuestran que son las más críticas para el
rendimiento del modelo. El \emph{logit lens} resuelve esta aparente paradoja:
las capas finales no \emph{crean} señal emocional, sino que la rotan
a la geometría del clasificador. La brecha entre probe F1 (información
presente) y gold sigmoid del \emph{logit lens} (información leíble)
cuantifica directamente este trabajo de alineación, que varía
sistemáticamente entre 0{,}63 (emociones tempranas) y 0{,}44 (tardías)
en su valor máximo. Las capas tardías son el cuello de botella entre representación y predicción, y por eso resultan frágiles bajo
compresión.

\paragraph{Las capas tardías hacen dos tipos de trabajo distintos.}
El análisis por grupo emocional revela que las capas L9--L11 no
realizan una única operación, sino dos cualitativamente distintas
según el contenido. Para información \emph{léxica ya formada}
(emociones tempranas como \emph{gratitude}), realizan
alineación pura: rotar una representación ya construida a la
base del clasificador. Para información \emph{contextual} que se
construye a lo largo del modelo (emociones tardías como
\emph{annoyance}), realizan formación y alineación coordinadas. La brecha entre probing y \emph{logit lens} distingue ambos
regímenes cuantitativamente.

\paragraph{Disociación entre geometría SVD y función.}
La compresión SVD opera sobre la estructura de valores singulares de
las matrices de pesos, una geometría que es ortogonal a la
organización funcional del modelo. La SVD no ``ve'' qué neuronas
codifican qué emociones: elimina señal uniformemente. Sin embargo, la
vulnerabilidad diferencial existe y se explica por la intensidad de
codificación: las emociones que el modelo escribe con más fuerza en
sus neuronas son más sensibles a cualquier perturbación de los pesos,
incluida la compresión.

Estas cuatro conclusiones fundamentan la estrategia de compresión
informada que se desarrolla en el Capítulo~\ref{sec:sintesis}, donde
la profundidad de cristalización y la intensidad de codificación
neuronal se utilizan como predictores para guiar decisiones de
compresión selectiva. El conjunto del capítulo establece que la
interpretabilidad mecánica no produce una colección de hallazgos
disjuntos sobre el modelo, sino una descripción \emph{unificada} de su
arquitectura funcional, descripción que cinco métodos
metodológicamente independientes corroboran.